\exnumber{04}
\extitle{UX III - Ex 5 - Personas}
\exfiles{ex04/Figma file and PDF, named Name\_FirstName\_DeliverableName\_Date\_VersionNumber}
\exauthorize{Those named in the exercise}
\makeheaderfiles

\textbf{Instructions:}

Let's create a persona based on your quantitative research and your interviews.

Your persona card should include common trends amongst the people you got data from (habits, verbatims, fears, drivers, frustration, and pain points).

Reminder: Your persona is fictional and should not have the exact name or picture as someone you interviewed, but remember a persona is here so you can make decisions as a designer.

\textbf{Deliverables:}

\begin{itemize}
\item A PDF persona card extracted from a Figma file
\item Upload the PDF document on Github, labeled as follows: "Name\_FirstName\_DeliverableName\_Date\_VersionNumber". For example, it can be labeled as "Doe\_John\_MidFiPrototype\_05242024\_V1".
\item Add a title within the document.
\item Pay attention to the global aspect of the document. It must be clear, simple and easy to read. You can get inspiration on canva.com but don't overload your design with too much details!
\end{itemize}

\textbf{Resources:}

\begin{itemize}
\item Make your persona great again \url{https://uxdesign.cc/personas-e60c1c06ead1}
\item Persona, journey maps, user flow, site map, OH MY! \url{https://uxdesign.cc/personas-journey-maps-site-maps-and-user-flows-oh-my-e71d044b4bcb}
\item Persona framework \url{https://miro.medium.com/max/1400/0*RbHfR6BHhhFIENNf}
\end{itemize}

\textbf{Quotes:}

\begin{itemize}
\item If you design for everyone, you delight no one.
\item A good user persona is the one based on user research, without regard to how many attributes we can describe.
\end{itemize}

